\subsection{High-dimensional Cox experiments: settings and numerical validation}
\label{app:cox-high-dimensional}

\paragraph{Datasets and motivation.}
These experiments test attribution at the dimensions produced by
genome-wide assays, where each patient contributes many molecular
measurements but the number of patients with survival follow-up is
comparatively small. GSE24080 contains multiple-myeloma expression
measurements, and TCGA lower-grade glioma (LGG) contains DNA methylation
measurements \cite{geo_gse24080,tcga_lgg_methylation}.
After filtering for valid overall-survival outcomes, the cohorts contain
553 and 511 patients, with 54,675 expression probe sets and 396,065
methylation sites, respectively. These dimensions count individual probes
or sites, without gene-level aggregation. We use the supplied 339/214
training/validation split for GSE24080 and a 383/128 training/test split
for TCGA LGG, stratified by event status with seed 42. Missing entries
are imputed using training medians, and centering and scaling are fitted
using training observations only.

\paragraph{Survival model and explanation target.}
We use Cox proportional-hazards models to incorporate patients whose
death has not been observed by the end of follow-up, called right-censored
observations. The model separates a common baseline hazard $h_0(t)$
from a patient-specific relative hazard:
\[
 h(t\mid x)=h_0(t)f(x),\qquad
 f(x)=\exp(\beta^\top x)=\prod_{j=1}^{d}\exp(\beta_jx_j).
\]
To stabilize estimation when $d\gg n$, a ridge penalty shrinks the
coefficients towards zero; the implementation's penalty parameter is
fixed in advance at $\alpha=0.01d$, without a hyperparameter search.
We use the Breslow approximation for patients with equal recorded event
times. All retained features have nonzero fitted coefficients.
The attribution target is the relative hazard $f(x)$, rather than the
log-relative hazard or a survival probability at a chosen time horizon.
Its feature-wise product structure is the property used by QuadraSHAP.

\paragraph{Background game and patient selection.}
To make comparisons refer to the same attribution problem, we fix 30
background patients sampled without replacement from the training set
using seed 42. For a feature subset $S$, we replace absent features
with entries from each background row $z^{(b)}$ and average the
resulting predictions:
\[
 v_x(S)=\frac{1}{30}\sum_{b=1}^{30}
 f\bigl(x_S,z^{(b)}_{\overline S}\bigr).
\]
Here, $\overline S$ denotes the absent features. This defines a uniformly
weighted sum of product games. We explain the first five held-out
patients in each saved split, without selection based on their outcomes.
The fitted coefficients, preprocessing, background rows, and patient
inputs are identical across methods. The quadrature guarantee concerns
this fixed empirical game; it does not bound the difference between a
30-row background and the population background distribution.

\paragraph{Quadrature rules and certification.}
We compare the Gauss--Legendre rule with
$m_{\mathrm{exact}}=\lceil d/2\rceil$ nodes against patient-specific
rules selected for $\varepsilon\in\{10^{-1},10^{-2},10^{-3}\}$.
The exact-degree rule integrates the relevant polynomial in real
arithmetic; it does not enumerate all $2^d$ coalitions or eliminate
floating-point rounding. For patient $p$, let $\phi_p$ denote the
Shapley vector of the fixed background game and $\phi_{m_p,p}$ the
quadrature result in real arithmetic. Node selection uses the analytical
bound $C_p$ to ensure
\[
 \|\phi_{m_p,p}-\phi_p\|_\infty\le C_p\le\varepsilon,
 \qquad C_{\max}=\max_{p=1,\ldots,5}C_p.
\]
The infinity norm measures the largest absolute error among features;
the tolerance is therefore absolute in the units of relative-hazard
attribution. Every selected rule satisfies this bound. The approximate
node counts span 44--54 for GSE24080 and 37--51 for TCGA LGG across
patients and tolerances, compared with 27,338 and 198,033 exact nodes.

\paragraph{Hardware, timing, and cached setup.}
We evaluate the GPU implementation on an Apple M4 Pro using JAX and
jaxlib 0.4.34 with jax-metal 0.1.1. The numerical kernel uses FP32,
with FP64 host accumulation, an 8\,GiB memory-planning budget, and one
background row per block.
An untimed call warms up each distinct dataset/node-count configuration.
We then record one synchronized explanation call per patient and method,
reporting the mean and sample standard deviation across the five
patients. The standard deviation describes variation across patients,
not uncertainty from repeated calls for one patient. Timings include
automatic node selection for approximate runs, data transfers, and host
accumulation; they exclude model fitting, warm-up, and quadrature-rule
construction. Approximation rows are fresh measurements; exact-degree
rows reuse the matched saved measurements after verifying the same
inputs, numerical source hashes, software versions, and memory plan.
Previously measured CPU construction of the cached exact rules took
10.97\,s and 561.92\,s, respectively. Fresh approximate-rule
construction took 0.386--0.563\,ms, reported as the range of medians
over five constructions per rule. Caching amortizes rule construction
across patients; each explanation still evaluates the quadrature
integrand for every background row.

\paragraph{Runtime and comparison with exact-degree GPU results.}
The mean exact-degree latencies are 1.172\,s for GSE24080 and
66.377\,s for TCGA LGG. Across the three tolerances, approximate
means range from 0.085--0.090\,s and 0.310--0.320\,s,
corresponding to approximately $13$--$14\times$ and
$207$--$214\times$ speedups relative to the reused exact means.
Table~\ref{tab:cox-high-dimensional} reports the largest discrepancy
between the floating-point approximation and its matching GPU
exact-degree output,
\[
 \Delta_{\mathrm{GPU}}=
 \max_p\|\widehat\phi^{\mathrm{GPU}}_{m_p,p}
 -\widehat\phi^{\mathrm{GPU}}_{\mathrm{exact},p}\|_\infty.
\]
For both $10^{-1}$ and $10^{-2}$, all five patients per dataset
meet their threshold under this comparison. At $10^{-3}$, only
4/5 GSE24080 and 2/5 TCGA patients do so; the largest discrepancies
are $1.31\times10^{-3}$ and $5.95\times10^{-3}$.

\paragraph{Independent precision checks.}
Agreement between two FP32 outputs can conceal numerical errors shared
by both computations. We therefore compare the GPU approximations with
a separate 64-node FP64 reference using numerically stable
\texttt{log1p}/\texttt{expm1} expressions. Its analytical quadrature
bounds are below $10^{-10}$, and a 96-node cross-check for the first
patient of each cohort changes the reference by less than $10^{-12}$.
This is a high-accuracy numerical reference, not a full degree-exact
FP64 evaluation. Relative to it, GPU pass counts for
$\varepsilon=(10^{-1},10^{-2},10^{-3})$ are $(5,5,1)$ out of five
for GSE24080 and $(5,3,0)$ for TCGA LGG. In particular, TCGA at
$10^{-2}$ passes for all patients against the GPU exact-degree
output, but only three against the independent reference, with a worst
error of $3.52\times10^{-2}$.

To assess the effect of arithmetic precision, we also evaluate the same
selected node counts and fixed inputs on CPU using FP64 and the same
parallel-reduction formula as the GPU kernel. All 30 patient--tolerance
cases meet their thresholds; the largest discrepancies from the
independent reference are $1.28\times10^{-11}$ for GSE24080 and
$1.59\times10^{-10}$ for TCGA LGG. These are separate accuracy
diagnostics and are not paired with the GPU timings. Their success
supports the adequacy of the selected quadrature rules in these cases,
while the FP32 comparisons expose a numerical accuracy limitation.
The analytical certificate controls quadrature approximation only;
these empirical precision checks do not extend it to a rigorous
floating-point error guarantee.
